\subsection{无穷乘积部分积的渐近展开}

我们知道 (Gen.~Top., V, p.~22 与 23)，以 \(1 + u_{n}\) （ \(u_{n} > -1\) ）为通项因子的无穷乘积收敛（相应地，交换收敛）的必要充分条件是：以 \(\log(1 + u_{n})\) 为通项的级数收敛（相应地，交换收敛），并且这时有关系

\[
\log \prod_ {n = 1} ^ {\infty} (1 + u _ {n}) = \sum_ {n = 1} ^ {\infty} \log (1 + u _ {n}).
\]

当无穷乘积收敛时，我们知道 \(u_{n}\) 趋于 0；于是 \(\log(1 + u_{n}) \sim u_{n}\) ；又我们知道，实数项级数交换收敛的必要充分条件是它绝对收敛 (Gen.~Top., IV, p.~372, prop. 5)；于是由命题 1 重新得到如下事实：以 \(1 + u_{n}\) 为通项因子的乘积交换收敛当且仅当以 \(u_{n}\) 为通项的级数绝对收敛 (Gen.~Top., IV, p.~368, th. 4)。

类似的论证适用于通项因子为复数 \(1 + u_{n}\) （ \(u_{n} \neq -1\) ）的无穷乘积。事实上，这样的乘积交换收敛的必要充分条件是 (Gen.~Top., VIII, p.~115, prop. 2)：以 \(|1 + u_{n}|\) 为通项因子的无穷乘积交换收敛，此外，若 \(\theta_{n}\) 是 \(1 + u_{n}\) 的辐角（取在 \(-\pi\) 与 \(+\pi\) 之间），则诸 \(\theta_{n}\) 所成的级数交换收敛。由于 \(u_{n}\) 趋于 0， \(\log(1 + u_{n})\) 从 n 的某个特定值起有定义 (III, p.~100)，且有

\[
\log (1 + u _ {n}) = \log | 1 + u _ {n} | + i \theta_ {n};
\]

于是，以 \(1 + u_{n}\) 为通项因子的乘积交换收敛的必要充分条件是：以 \(|\log(1 + u_{n})|\) 为通项的级数绝对收敛 (Gen.~Top., VII, p.~84, th. 1)；而 \(\log(1 + u_{n}) \sim u_{n}\) (I, p.~18, prop. 5)，故重新得到以 \(u_{n}\) 为通项的级数绝对收敛这一条件 (Gen.~Top., VIII, p.~116, th. 1)。

无穷乘积与实数项级数之间的关系有时使我们能得到部分积 \(p_{n} = \prod_{k=1}^{n}(1 + u_{k})\) 的渐近展开；只需有部分和 \(s_{n} = \sum_{k=1}^{n} \log(1 + u_{k})\) 的渐近展开，再展开 \(p_{n} = \exp(s_{n})\) 即可；于是问题归结为前面考察过的两个问题 (V, p.~238 与 p.~226)。

例：斯特林公式. 我们来求 \(n!\) 的渐近展开；这导致我们展开 \(s_n = \sum_{p=1}^{n} \log p\) ，然后展开 \(\exp(s_n)\) 。 \(n^{\circ} 2\) 的方法依次给出

\[
s _ {n} = \sum_ {p = 1} ^ {n} \log p \sim \int_ {1} ^ {n} \log t d t = n \log n - n + 1
\]

继而

\[
\log n - \int_ {n - 1} ^ {n} \log t d t = \log n - (n \log n - (n - 1) \log (n - 1) - 1) \sim \frac {1}{2 n}
\]

由此

\[
s _ {n} = n \log n - n + \frac {1}{2} \log n + o (\log n).
\]

然后

\[
\log n - \int_ {n - 1} ^ {n} \log t d t - \frac {1}{2} (\log n - \log (n - 1)) \sim - \frac {1}{1 2 n ^ {2}}
\]

由此

\[
s _ {n} = n \log n - n + \frac {1}{2} \log n + k + \frac {1}{1 2 n} + o _ {1} \left(\frac {1}{n}\right)
\]

（k 为常数）

最后由此推出 (V, p.~226)

\[
n! = e ^ {k} n ^ {n + 1 / 2} e ^ {- n} \left(1 + \frac {1}{1 2 n} + o _ {2} \left(\left(\frac {1}{n}\right)\right)\right).\tag{5}
\]

我们将在 VII, p.~322 证明 \(e^k = \sqrt{2\pi}\) 。公式 (5)（取 \(k\) 的这个值）称为斯特林公式。同样，对每个实数 \(a\) （非 \(>0\) 的整数），可以证明

\[
(a + 1) (a + 2) \dots (a + n) \sim \mathrm{K} (a) n ^ {n + a + \frac {1}{2}} e ^ {- n}.\tag{6}
\]

我们也将在 (VII, p.~18) 确定函数 \(\mathrm{K}(a)\) 。由公式 (5) 与 (6) 特别可推出

\[
\binom {a} {n} \sim (- 1) ^ {n}   \varphi (a)   n ^ {- a - 1}\tag{7}
\]

对每个实数 \(a\) （非 \(> 0\) 的整数）成立，其中 \(\varphi(a)\) 是 \(a\) 的一个函数，将在 VII, p.~322 明确给出。

